\chapter{Repository-Aware Context Retrieval}
\label{cha:ch1}

\section{Existing Context Retrieval Algorithms}
% \myx{This section should not call "Introduction" (there can be only one intro in the whole paper). Maybe a better section title is "Existing Repository-Aware Context Retrieval Algorithms"?}

Modern code LLM agents have shown impressive capabilities, but most still struggle with incomplete or irrelevant context when working on large repositories. Many frameworks, from early tool-using agents to recent systems like SWE-Agent \cite{yang2024sweagent}, AutoCodeRover \cite{zhang2024autocoderover}, and Agentless \cite{xia2024agentlessdemystifyingllmbasedsoftware}, often feed the model only limited snippets of a codebase or generic summaries. Although agents often are able to interact with the entire codebase via some "scrolling" mechanism \myx{would be good to add a brief description} \mytodo{add SWE-agent mechanism as example here}, this does not prevent them from getting stuck in a local optimum \cite{ouyang2025repographenhancingaisoftware} \myx{what does "local optimum" refer to, better be specific}. This lack of essential context can lead to unstable performance and failures in code generation or bug-fixing (we refer to this as \textbf{Challenge 1}). For example, providing an agent with only a function implementation but not definitions of its referenced symbols (classes, functions, etc.) can omit crucial information and hurt the quality of patches\myx{A figure to illustrate this example would be good.}. Conversely, naively supplying too much code (e.g., multiple flat files or irrelevant documentation/config) can overwhelm the model with irrelevant details, interfering with its reasoning \myx{Would be good to add some citation here.}. Previous studies have noted the challenge of balancing concise yet sufficient context given the finite context window of LLMs \cite{ouyang2025repographenhancingaisoftware}.

% Explorations have been made in an effort to mitigate Challenge 1. A common approach is retrieval-augmented generation (RAG), where the agent searches the repository for relevant snippets given an issue description. For instance, some systems treat the bug report as a query and use information retrieval (e.g.\ BM25 or dense embeddings) to pull in code segments likely related to the issue. Sliding window scans have also been used to ensure no part of a large file is missed by the context. Other frameworks rely on the LLM itself to navigate code: Agentless presents the entire repository structure (files and directories) to the model, letting it decide which files to open and read in a hierarchical fashion \cite{xia2024agentlessdemystifyingllmbasedsoftware}. In contrast, graph-based indexing methods pre-compute a knowledge graph of the repository (linking function definitions to references, etc.) and retrieve context via graph queries. These can provide stronger guidance, but maintaining a full repository graph (or vector index) is resource-intensive and struggles with continuously evolving codebases.

RepoGraph offers a state-of-the-art solution to Challenge 1 by explicitly modeling the structure and dependencies of the repository. It constructs a comprehensive \emph{repository-level code graph} that captures relationships between code elements across different files \cite{ouyang2025repographenhancingaisoftware}. Each class, method, or function is a node, and the edges connect definitions to their usage or dependency. This enables the agent to gather all relevant code context needed for a fix without having to traverse the entire codebase. \mytodo{include a graphic example here}. RepoGraph leverages repository-wide structure to provide full and fine-grained context. By following these links, the agent is much less likely to overlook a relevant piece of code compared to linear search or limited file context. This approach has proven effective with GPT4/4o or Claude-3.5-Sonnet as backend: when plugged into Agentless or SWE-Agent, RepoGraph consistently boosted their success rates \cite{ouyang2025repographenhancingaisoftware}. This makes RepoGraph ideal as the foundation of our frontend context retrieval algorithm, whose specific usage will be discussed in closer detail in Chapter \ref{cha:ch3}.

\mytodo{eval related insights, especially why comprehensive eval with diff APIs is important}